\documentclass[11pt]{article}

% ---------------------------------------------------------------------------
% Spoken teaching notes for a ~2-hour hands-on practice session on
% Diffusion Models and Diffusion Policy (5 Jupyter notebooks).
% Author: Alessandro Leanza (SUPSI)
% Compiles with a plain: pdflatex practice_notes.tex
% ---------------------------------------------------------------------------

\usepackage[utf8]{inputenc}
\usepackage[T1]{fontenc}
\usepackage{amsmath}
\usepackage{amssymb}
\usepackage[margin=2.5cm]{geometry}
\usepackage{xcolor}
\usepackage{booktabs}
\usepackage{longtable}
\usepackage{enumitem}
\usepackage{listings}
\usepackage{hyperref}

\hypersetup{
  colorlinks=true,
  linkcolor=blue!55!black,
  urlcolor=blue!55!black,
  citecolor=blue!55!black,
  pdftitle={Hands-on Practice: Diffusion Models and Diffusion Policy},
  pdfauthor={Alessandro Leanza}
}

% --- Colours ---------------------------------------------------------------
\definecolor{saybg}{RGB}{232,242,253}
\definecolor{sayrule}{RGB}{33,118,209}
\definecolor{codebg}{RGB}{247,247,247}
\definecolor{codekw}{RGB}{0,90,160}
\definecolor{codecomment}{RGB}{110,140,110}
\definecolor{codestr}{RGB}{170,70,40}
\definecolor{hlcolor}{RGB}{150,30,30}

% --- "Say this" environment (robust: built on quote + colour) --------------
% A coloured, left-ruled paragraph block. No mdframed / tcolorbox needed.
\newenvironment{saythis}{%
  \par\smallskip\noindent
  \begin{list}{}{%
    \setlength{\leftmargin}{1.2em}%
    \setlength{\rightmargin}{0.5em}%
    \setlength{\topsep}{2pt}%
  }\item\relax
  \color{sayrule}\textbf{Say this.}\hspace{0.4em}\color{black}\itshape
}{%
  \end{list}\par\smallskip
}

% --- Highlight / ask environment -------------------------------------------
\newcommand{\highlight}[1]{\par\smallskip\noindent\textcolor{hlcolor}{\textbf{Highlight / ask.}}\hspace{0.4em}#1\par\smallskip}

% --- "What it does" tag ----------------------------------------------------
\newcommand{\whatdoes}[1]{\par\noindent\textbf{What it does.}\hspace{0.4em}#1\par}

% --- Listings configuration for Python -------------------------------------
\lstset{
  language=Python,
  basicstyle=\ttfamily\small,
  backgroundcolor=\color{codebg},
  keywordstyle=\color{codekw}\bfseries,
  commentstyle=\color{codecomment}\itshape,
  stringstyle=\color{codestr},
  showstringspaces=false,
  breaklines=true,
  frame=single,
  framerule=0.3pt,
  rulecolor=\color{gray!40},
  xleftmargin=0.5em,
  xrightmargin=0.5em,
  columns=fullflexible,
  keepspaces=true,
  aboveskip=6pt,
  belowskip=6pt
}

\setlist[itemize]{leftmargin=1.4em,topsep=2pt,itemsep=1pt}
\setlength{\parskip}{3pt}
\setlength{\parindent}{0pt}

\title{\textbf{Hands-on Practice: Diffusion Models and Diffusion Policy}\\[0.3em]
\large Spoken Teaching Notes for a 2-Hour Session (5 Notebooks)}
\author{Alessandro Leanza \\ \small SUPSI}
\date{}

\begin{document}
\maketitle

\begin{abstract}
\noindent
These are the instructor's spoken notes for a two-hour hands-on practice session that
accompanies the lecture on Diffusion Models and Diffusion Policy for compliant robots.
The session works through five Jupyter notebooks, run on Google Colab with a GPU.
The unifying theme is that the \emph{same} diffusion process operates on images, video, and
robot trajectories --- only the neural-network backbone changes (a 2D U-Net for images, a 1D
U-Net or Transformer for sequences). Notebook~01 builds data intuition; notebook~02 visualises
the forward and reverse processes; notebook~03 trains a full DDPM on MNIST from scratch;
notebook~04 benchmarks four backbones; notebook~05 trains a real Diffusion Policy on the Push-T
dataset and shows why fast sampling (DDIM) matters for real-time control. For every code block
below the notes give three things: \textbf{What it does}, a ready-to-read \textbf{Say this} script,
and \textbf{Highlight / ask} points (with a question to pose to students). The two long training
runs (notebooks~03 and~05) are launched early and used as cover for the theory, so the wall-clock
fits into 120 minutes.
\end{abstract}

\tableofcontents
\bigskip

% ===========================================================================
\section{Session Plan (2 hours)}
% ===========================================================================

The single most important wall-clock fact: \textbf{notebooks 01, 02 and 04 run in seconds}
(synthetic data and forward passes only), while \textbf{notebook 03 trains for $\approx$15--30~min}
(DDPM on MNIST, \texttt{NUM\_EPOCHS = 20}) and \textbf{notebook 05 trains for $\approx$30--40~min}
(Diffusion Policy on Push-T, \texttt{NUM\_EPOCHS = 100}, a checkpoint saved every 10 epochs) on a
Colab T4. The strategy is to \emph{overlap} those two trainings with explanation.

\paragraph{Recommended ordering.} Run notebooks 01 and 02 quickly (they are pure visualisation).
Then open notebook 03, walk through its setup and model, and \textbf{start its training cell}.
While the MNIST DDPM trains, cover all of notebook 04 (architectures --- fast forward passes) and
the DDPM theory. By the time you finish 04, notebook 03 is done, so you show its samples. Then
open notebook 05, run setup and the dataset download, build the model, and \textbf{start its
training cell}; lecture through the Diffusion Policy theory and the lecture bridge while it trains,
then return to show the sampled action chunks and the DDPM-vs-DDIM speed comparison.

\begin{center}
\begin{longtable}{@{}p{2.0cm} p{4.6cm} p{2.0cm} p{6.0cm}@{}}
\toprule
\textbf{Time} & \textbf{Activity} & \textbf{Wall clock} & \textbf{Notes} \\
\midrule
\endhead
0--10 min   & NB 01 --- Data modalities          & runs in seconds & Images, video, trajectories; normalization. \\
10--22 min  & NB 02 --- Forward / reverse        & runs in seconds & Schedules, SNR, oracle DDPM \& DDIM reverse. \\
22--30 min  & NB 03 --- Setup + model walk-through & seconds        & Cover dataset, schedule, U-Net; \textbf{then launch training}. \\
30--32 min  & \textbf{Start NB 03 training}      & 15--30 min (bg) & Click run on the training cell, let it run. \\
32--55 min  & NB 04 --- Architectures (all)      & runs in seconds & 2D U-Net, +attn, 1D U-Net, DiT; benchmark. Talk over training. \\
55--62 min  & Back to NB 03 --- results          & seconds         & Loss curve, samples, denoising trajectory. \\
62--75 min  & NB 05 --- Setup, download, data    & $\sim$1 min dl  & Zarr load, feature engineering, dataset windows, model. \\
75--77 min  & \textbf{Start NB 05 training}      & 30--40 min (bg) & Launch training cell. \\
77--108 min & Diffusion Policy theory + lecture bridge & ---       & Conditioning, action chunking, DDPM$\to$DDIM, real-time control. \\
108--118 min& Back to NB 05 --- results          & seconds         & Sampled action chunks, DDPM vs DDIM speed, distributions. \\
118--120 min& Wrap-up / Q\&A                     & ---             & Key takeaways and the bridge to the lecture. \\
\bottomrule
\end{longtable}
\end{center}

\highlight{If the GPU is slow or time is short, both trainings can be cut: notebook 03 still
produces recognizable digits after $\approx$10 epochs, and notebook 05 saves a usable checkpoint
every 10 epochs. The notebook 05 final cell can also \emph{load a checkpoint} to skip training
entirely. Always start the two long cells \emph{before} you start lecturing.}

\paragraph{A note on the Setup cells.} CODE \#1 in every notebook is an added ``Setup'' cell:
it imports \texttt{torch} and prints the PyTorch version and whether CUDA is available. In notebook
05 it also \texttt{pip install}s \texttt{zarr} and \texttt{gdown} if missing. Run it first in each
notebook and confirm the line says \texttt{CUDA available: True}; if it says \texttt{False}, fix the
runtime (\texttt{Runtime $\to$ Change runtime type $\to$ GPU}) before going further. I will not
re-explain this cell for every notebook.

% ===========================================================================
\section{Notebook 01 --- Data Modalities}
% ===========================================================================

\textbf{Learning goal.} Build the intuition that a diffusion model just sees \emph{tensors}: it does
not care whether a sample is an image, a video clip, or a robot trajectory. We look at the PyTorch
shapes of each modality, count their dimensionality $d$, and see why every modality must be
normalized to roughly $\mathcal{N}(\mathbf{0}, \mathbf{I})$. This is the conceptual foundation for
the whole session.

\subsection{Setup (CODE \#1) and library imports (CODE \#2)}
\whatdoes{CODE \#1 checks PyTorch / CUDA. CODE \#2 imports NumPy, Matplotlib, PyTorch and PIL,
sets a seaborn style, and fixes the random seeds (\texttt{np.random.seed(42)},
\texttt{torch.manual\_seed(42)}) for reproducibility.}
\begin{saythis}
Before we touch any maths, let's just load our tools. We fix the random seed so that everything you
see on my screen you can reproduce exactly on yours. From now on, every number in this notebook is
deterministic.
\end{saythis}
\highlight{Seeds make the demos reproducible. Ask: \emph{why might we want a fixed seed in a teaching
demo but a random seed in real research?}}

\subsection{Synthetic RGB image and channel decomposition (CODE \#3)}
\whatdoes{Builds a $64\times64\times3$ synthetic image (a red/green gradient, constant blue, a yellow
circle, a white corner square) and plots the RGB image plus its three separate channels. The array is
in \texttt{(H, W, C)} ``channel-last'' NumPy layout with values in $[0,1]$.}
\begin{saythis}
An image is nothing more than a grid of numbers. Here I have made a synthetic one so we don't depend
on any file. Notice the layout: height, then width, then the three colour channels. Each channel is
itself just a 2D array of intensities --- red, green, blue --- and stacking them gives the colour
image you see on the left.
\end{saythis}
\highlight{Stress the NumPy/PIL convention \texttt{(H, W, C)} versus the PyTorch convention
\texttt{(C, H, W)} --- this trips up everyone at some point. Ask: \emph{how many numbers does a single
$64\times64$ RGB image contain?} (Answer: $3\times64\times64 = 12{,}288$.)}

\subsection{Tensor format and batching (CODE \#4)}
\whatdoes{Permutes the image to PyTorch \texttt{(C, H, W)}, prints the total element count
$d = C\!\times\!H\!\times\!W$, then builds a batch of 4 by repeating the image with small brightness
offsets to get \texttt{(B, C, H, W)}. Finally it normalizes the batch from $[0,1]$ to $[-1,1]$.}
\begin{saythis}
Neural networks in PyTorch want channels first, so we permute to ``C, H, W''. And they never see one
image at a time --- they see a \emph{batch}. So the real shape the model works with is
``B, C, H, W''. The last line is the one I want you to remember: we rescale pixels from zero-to-one to
minus-one-to-one. That centring around zero is exactly what the diffusion noise schedule expects.
\end{saythis}
\highlight{The $[-1,1]$ normalization is not cosmetic: the forward process ends at
$\mathcal{N}(\mathbf{0},\mathbf{I})$, so zero-mean data makes the signal-to-noise schedule clean. We
will see this again in every notebook.}

\subsection{Synthetic video (CODE \#5)}
\whatdoes{Generates an 8-frame $32\times32$ video of an orange ball moving left to right on a colour-shifting
background, shape \texttt{(T, H, W, C)}; permutes to \texttt{(T, C, H, W)}, adds a batch dimension to get
\texttt{(B, T, C, H, W)}, and shows all 8 frames.}
\begin{saythis}
A video is just a stack of images along a new axis: time. So we add a $T$ dimension at the front and
get ``T, C, H, W'' for one clip, or ``B, T, C, H, W'' for a batch. Conceptually nothing has changed ---
it is still a tensor of numbers --- but look at how the total dimensionality explodes once you add
time. This is exactly why video models like Sora work in a compressed latent space rather than on raw
pixels.
\end{saythis}
\highlight{Dimensionality grows multiplicatively with $T$. This motivates latent diffusion (LDM).
Ask: \emph{if one frame is $12{,}288$ numbers, how big is a 2-second clip at 30 fps?}}

\subsection{Robot trajectory datasets (CODE \#6)}
\whatdoes{Defines a figure-eight 2D end-effector trajectory and a 6-DOF joint trajectory, both of length
$T=100$; builds a dataset of 20 noisy figure-eight ``demonstrations'' to mimic an imitation-learning
dataset. Shapes are \texttt{(T, D)} per trajectory and \texttt{(B, T, D)} for the batch.}
\begin{saythis}
Now the punchline of this notebook. A robot trajectory is \emph{also} just a tensor. Here each
demonstration is a sequence of two-dimensional end-effector positions --- a figure-eight ---
with a little noise, like 20 slightly different human demonstrations. The shape is ``T by D'':
sequence length by action dimension. This ``B, T, D'' tensor is precisely what Diffusion Policy will
generate in notebook five.
\end{saythis}
\highlight{This is the conceptual hinge of the day: trajectories live in the same tensor world as
images. The action dimension $D$ is small (2, 3, 6, 7), but the temporal structure is what matters.}

\subsection{Trajectory and dimensionality plots (CODE \#7)}
\whatdoes{Plots the 2D trajectory dataset, the 6 joint-angle curves over time, and a bar chart comparing
the total dimensionality $d$ across an MNIST image (784), an RGB image (12{,}288), a video clip (24{,}576),
a joint trajectory (600) and a 2D trajectory (200).}
\begin{saythis}
Two things to read off these plots. First, the demonstrations form a \emph{distribution} of
trajectories, not a single path --- and learning that distribution is exactly what a generative model
does. Second, look at the bar chart: a robot trajectory is tiny compared with a video. That is good
news --- diffusion on trajectories is cheap enough to run on a robot in real time, which is the whole
point of Diffusion Policy.
\end{saythis}
\highlight{Trajectories are low-dimensional, which is why real-time diffusion control is feasible.
Contrast this with video. Ask: \emph{which of these modalities would benefit most from a latent space?}}

\subsection{Normalization examples (CODE \#8)}
\whatdoes{Normalizes the image from $[0,1]$ to $[-1,1]$ and standardizes the trajectory batch per
dimension, $(x-\mu)/\sigma$. Prints before/after statistics and plots pixel and trajectory value
histograms, raw versus normalized.}
\begin{saythis}
Different modalities need different normalization recipes, but they all share one goal: get the data
roughly zero-mean and unit-variance. For images we use a fixed linear map to minus-one-to-one. For
trajectories we standardize each dimension using the dataset's own mean and standard deviation. Watch
the histograms shift to be centred on zero --- that is the distribution the diffusion prior assumes.
\end{saythis}
\highlight{Per-dimension standardization for trajectories is exactly what notebook 05 does to
observations and actions, and the inverse map is applied at sampling time. If you forget to
de-normalize the generated actions, the robot moves nonsensically.}

% ===========================================================================
\section{Notebook 02 --- Forward and Reverse Diffusion}
% ===========================================================================

\textbf{Learning goal.} Make the two halves of DDPM concrete and visual: the \emph{forward} process
$q$ that destroys an image into noise, and the \emph{reverse} process $p_\theta$ that rebuilds it.
We compare linear and cosine noise schedules, read off the signal-to-noise ratio, and run an
\emph{oracle} reverse process (using the true stored noise) for both DDPM and DDIM to see what a
perfect denoiser would achieve.

\subsection{Setup, imports, configuration (CODE \#1--3)}
\whatdoes{CODE \#1 checks CUDA. CODE \#2 imports and seeds. CODE \#3 is the configuration cell:
\texttt{IMAGE\_PATH = None} (use synthetic), \texttt{IMAGE\_SIZE = 64}, \texttt{T\_STEPS = 1000} total
diffusion steps, and the device.}
\begin{saythis}
The configuration cell is where you would plug in your own photo if you wanted to. We will use the
synthetic image and the standard one-thousand-step schedule from the original DDPM paper.
\end{saythis}
\highlight{$T = 1000$ is the canonical DDPM choice. We will see notebook 05 use only $T = 100$ ---
small action chunks need far fewer steps.}

\subsection{Image utilities and the clean image (CODE \#4)}
\whatdoes{Defines helpers to make the synthetic image, load an optional file, and convert a $[-1,1]$
tensor back to a displayable $[0,1]$ image. Loads $\mathbf{x}_0$ as a \texttt{(1, C, H, W)} tensor in
$[-1,1]$ and displays it.}
\begin{saythis}
This is our clean starting point, $\mathbf{x}_0$. Everything from here is about corrupting this image
and then trying to recover it. Note it is already normalized to minus-one-to-one, as we discussed.
\end{saythis}

\subsection{Noise schedules: linear vs cosine (CODE \#5)}
\whatdoes{Implements the linear schedule ($\beta_t$ from $10^{-4}$ to $0.02$) and the cosine schedule
(Nichol \& Dhariwal). Computes $\alpha_t = 1-\beta_t$, the cumulative product
$\bar{\alpha}_t = \prod_{s\le t}\alpha_s$, and the mixing coefficients $\sqrt{\bar\alpha_t}$ and
$\sqrt{1-\bar\alpha_t}$. Plots $\beta_t$, $\bar\alpha_t$, and the coefficients for both schedules.}
\begin{saythis}
The schedule decides \emph{how fast} we destroy the signal. Beta is the variance of noise added at each
step. Alpha-bar, the cumulative product, is the fraction of the original signal that survives to step
$t$ --- so the middle plot, alpha-bar going from one down to zero, is literally the image dissolving.
The linear schedule destroys information very fast early on; the cosine schedule, in red, is gentler at
the start and keeps fine detail alive longer. That is why the cosine schedule trains better, and it is
the one we use for the rest of the day.
\end{saythis}
\highlight{Key relation to put on the board:
$\mathbf{x}_t = \sqrt{\bar\alpha_t}\,\mathbf{x}_0 + \sqrt{1-\bar\alpha_t}\,\boldsymbol\varepsilon$.
The signal coefficient and the noise coefficient always satisfy
$(\sqrt{\bar\alpha_t})^2 + (\sqrt{1-\bar\alpha_t})^2 = 1$. Ask: \emph{at what $t$ is half the signal
gone?}}

\subsection{Forward process visualization (CODE \#6)}
\whatdoes{Defines \texttt{q\_sample} implementing the closed-form reparameterization
$\mathbf{x}_t = \sqrt{\bar\alpha_t}\mathbf{x}_0 + \sqrt{1-\bar\alpha_t}\boldsymbol\varepsilon$.
Uses a \emph{single fixed} noise instance across all timesteps and shows the noisy image at 12
timesteps (top row) and the scaled noise component (bottom row), with the SNR printed per panel.}
\begin{saythis}
Here is the trick that makes diffusion trainable. We do \emph{not} add noise one step at a time. The
reparameterization lets us jump straight to any timestep $t$ in one line: a bit of the clean image
plus a bit of Gaussian noise, with the mixing set by the schedule. Watch the top row: at $t=0$ it is
the clean image, and by $t=999$ it is essentially pure noise. I deliberately reuse the \emph{same}
noise sample for every column, so you can see it is the \emph{coefficients} changing, not the noise.
\end{saythis}
\begin{lstlisting}
def q_sample(x0, t_idx, sched, noise=None):
    s = sched['sqrt_a_bar'][t_idx].view(-1,1,1,1)
    n = sched['sqrt_one_minus'][t_idx].view(-1,1,1,1)
    return s * x0 + n * noise, noise   # x_t and the epsilon used
\end{lstlisting}
\highlight{Closed-form sampling at any $t$ is \emph{the} reason DDPM training is efficient --- no need
to chain $t$ sequential steps per gradient update. The function also returns the noise $\varepsilon$,
because that noise is the training target.}

\subsection{Signal-to-noise ratio (CODE \#7)}
\whatdoes{Computes SNR in dB, $20\log_{10}(\sqrt{\bar\alpha_t}/\sqrt{1-\bar\alpha_t})$, for both
schedules; plots them, shades signal-dominant versus noise-dominant regions, and prints the timestep
where SNR crosses 0~dB for each schedule.}
\begin{saythis}
The signal-to-noise ratio gives a single number for ``how much real image is left''. Where the curve
crosses zero decibels, signal and noise are equal --- that is the tipping point. Notice the cosine
schedule stays in the signal-dominant region much longer. The denoiser spends most of its capacity
near that crossover, where the task is hardest.
\end{saythis}
\highlight{SNR connects the schedule to learning difficulty: the network learns most around the
$0$~dB region. Ask: \emph{which schedule gives the model more ``easy'' high-SNR examples to learn fine
detail?}}

\subsection{DDPM oracle reverse step and trajectory (CODE \#8--9)}
\whatdoes{CODE \#8 precomputes the posterior variance $\tilde\beta_t$ and defines one DDPM reverse step
using the \emph{true} stored noise as if it were the network prediction (the ``oracle''):
$\boldsymbol\mu_\theta = \frac{1}{\sqrt{\alpha_t}}\big(\mathbf{x}_t - \frac{\beta_t}{\sqrt{1-\bar\alpha_t}}\boldsymbol\varepsilon\big)$,
plus posterior-variance noise when $t>0$. CODE \#9 starts at $\mathbf{x}_T$ and iterates all the way
back to $\mathbf{x}_0$, visualizing $\sim$12 intermediate steps.}
\begin{saythis}
Now we run the process backwards. In a real model a neural network would predict the noise; here I
\emph{cheat} and feed in the true noise we actually added. This is the ``oracle'' --- it shows you the
best case, what a perfect denoiser could ever do. Starting from pure static, each step removes a
little noise according to the DDPM update, and the image re-emerges. Importantly, DDPM is
\emph{stochastic}: at every step except the last we inject a fresh bit of Gaussian noise, scaled by the
posterior variance.
\end{saythis}
\highlight{The oracle separates two questions: ``is the reverse \emph{math} correct?'' (yes, the image
comes back) versus ``can a network \emph{learn} the noise?'' (that's notebook 03). Emphasize the
stochasticity term --- it is what distinguishes DDPM from DDIM next.}

\subsection{DDIM oracle reverse step and trajectory (CODE \#10--11)}
\whatdoes{Defines the DDIM reverse step. It first predicts the clean image
$\hat{\mathbf{x}}_0 = (\mathbf{x}_t - \sqrt{1-\bar\alpha_t}\,\boldsymbol\varepsilon)/\sqrt{\bar\alpha_t}$,
then steps to $\mathbf{x}_{t-1} = \sqrt{\bar\alpha_{t-1}}\,\hat{\mathbf{x}}_0 + \sqrt{1-\bar\alpha_{t-1}-\sigma_t^2}\,\boldsymbol\varepsilon + \sigma_t \mathbf{z}$,
where $\eta$ controls stochasticity. With \texttt{eta = 0.0} the update is fully deterministic. CODE \#11
runs the oracle DDIM trajectory.}
\begin{saythis}
DDIM is the same forward process but a different \emph{reverse} rule. Instead of sampling from a
Gaussian posterior, it first estimates the clean image directly, then re-noises it to the previous
level. The parameter eta controls randomness: eta equals one behaves like DDPM, eta equals zero is
fully deterministic --- no random term at all. Deterministic reverse is the key idea, because it lets
us \emph{skip steps} and sample in far fewer iterations. Hold that thought: in notebook five it becomes
the difference between a robot that can react in real time and one that cannot.
\end{saythis}
\begin{lstlisting}
# DDIM: predict clean x0, then jump to the previous noise level
x0_pred = (x_t - torch.sqrt(1 - a_bar_t) * eps) / torch.sqrt(a_bar_t)
x_prev  = torch.sqrt(a_bar_prev) * x0_pred + dir_xt + sigma_t * z
\end{lstlisting}
\highlight{Same training, two samplers. DDIM with $\eta=0$ is deterministic and supports large step
strides $\Rightarrow$ fast sampling. Ask: \emph{why can a deterministic sampler safely take bigger
steps than a stochastic one?}}

\subsection{Forward vs reverse side-by-side (CODE \#12)}
\whatdoes{Lays out three rows over matched timesteps: forward $q$ (top), reverse DDPM (middle), reverse
DDIM (bottom), so you can compare the destruction and the two reconstructions step by step.}
\begin{saythis}
This single figure is the summary of the notebook. Top row: the forward process eating the image.
Middle and bottom: the two reverse processes rebuilding it. Read it left-to-right on top and notice the
reverse rows put it back together. With the oracle noise, both DDPM and DDIM recover the image
faithfully --- the difference between them is speed and determinism, which we will exploit later.
\end{saythis}
\highlight{Takeaway table for the board: forward closed-form = jump to any $t$; reverse = iterate;
network learns $\varepsilon_\theta \approx \varepsilon$; loss is plain MSE. This sets up notebook 03.}

% ===========================================================================
\section{Notebook 03 --- Complete DDPM Training Pipeline}
% ===========================================================================

\textbf{Learning goal.} Build a full Denoising Diffusion Probabilistic Model from scratch and
\emph{train} it on MNIST: cosine schedule, a U-Net with sinusoidal time embeddings, AdaGN time
conditioning and a bottleneck self-attention block, the simple MSE noise-prediction loss
$\mathcal{L}_{\text{simple}}$, and a 1000-step ancestral sampler. \textbf{Strategy: launch the training
cell early (CODE \#9) and lecture through notebook 04 while it runs.}

\subsection{Setup, imports, configuration (CODE \#1--3)}
\whatdoes{CODE \#1 checks CUDA; CODE \#2 imports torch / torchvision / tqdm and seeds. CODE \#3 is the
config: \texttt{IMG\_SIZE = 28}, \texttt{CHANNELS = 1}, \texttt{T\_STEPS = 1000}, \texttt{BATCH\_SIZE = 128},
\texttt{LR = 2e-4}, \textbf{\texttt{NUM\_EPOCHS = 20}}, \texttt{MODEL\_DIM = 64}.}
\begin{saythis}
Here is the whole experiment in one cell. Twenty epochs on MNIST will give us recognizable digits in
roughly fifteen to thirty minutes on this T4 GPU. That is too long to watch in silence, so the plan is:
we will read through the model now, hit ``run'' on the training cell, and then jump to notebook four
while this trains in the background.
\end{saythis}
\highlight{Point at \texttt{NUM\_EPOCHS = 20}. Publication quality needs 100+; 20 is the teaching
compromise. The comment in the code even says so.}

\subsection{MNIST dataset (CODE \#4)}
\whatdoes{Builds a torchvision MNIST \texttt{DataLoader} with \texttt{Normalize((0.5,),(0.5,))} to map
pixels to $[-1,1]$, prints sizes, and shows a $4\times4$ grid of training digits.}
\begin{saythis}
MNIST: sixty thousand twenty-eight by twenty-eight grayscale digits. The only preprocessing that matters
is the normalize step --- same idea as notebook one, mapping pixels to minus-one-to-one so the data
matches the noise prior.
\end{saythis}
\highlight{Same normalization principle as NB01, now in a real \texttt{transforms} pipeline. Ask:
\emph{what would happen to training if we left the data in $[0,1]$?}}

\subsection{Cosine noise schedule (CODE \#5)}
\whatdoes{Defines the cosine $\beta$ schedule and precomputes \texttt{alphas\_cumprod}, its previous-step
shift, $\sqrt{\bar\alpha_t}$, $\sqrt{1-\bar\alpha_t}$, $\sqrt{1/\alpha_t}$, and the posterior variance.
Defines \texttt{extract} (gathers per-sample schedule values) and plots $\bar\alpha_t$ and the mixing
coefficients.}
\begin{saythis}
Exactly the cosine schedule from notebook two, but now we precompute every quantity we will need for
training and sampling and cache it on the GPU. The little \texttt{extract} helper just picks out the
right schedule value for each example in the batch, given that each example gets a different random
timestep.
\end{saythis}
\highlight{Precomputing schedule tensors once is a standard efficiency pattern --- do it on the device,
index with \texttt{extract}.}

\subsection{Forward process on MNIST (CODE \#6)}
\whatdoes{Defines \texttt{q\_sample} (closed-form $\mathbf{x}_t$) and shows 4 digits noised at 6 timesteps,
visually confirming the forward process on real data.}
\begin{saythis}
Same closed-form forward sampling as before, now on actual digits. Notice a ``4'' is still faintly
visible at $t$ equals two-fifty, mostly gone by five hundred, and pure noise by the end. This is the
supervision signal: at training time we will pick a random $t$, noise the digit this much, and ask the
network what noise we added.
\end{saythis}

\subsection{U-Net building blocks (CODE \#7)}
\whatdoes{Defines (a) \texttt{SinusoidalEmbedding} encoding the integer timestep with sine/cosine
frequencies; (b) \texttt{ResBlock} with two conv layers and \emph{AdaGN} time conditioning --- the time
embedding produces a scale and shift applied after group norm; (c) a bottleneck \texttt{SelfAttention}
block.}
\begin{saythis}
Three ideas power the U-Net. First, the network must know \emph{which} timestep it is denoising, so we
encode $t$ with the same sinusoidal trick used for positions in Transformers. Second, that time
embedding is injected into every residual block as a learned scale-and-shift after group norm --- this
is AdaGN, adaptive group normalization. It tells each block ``you are at a high-noise step'' or ``a
low-noise step'' and the block adapts. Third, one self-attention block at the bottleneck lets distant
pixels talk to each other, which convolutions alone cannot do.
\end{saythis}
\begin{lstlisting}
# AdaGN time conditioning inside a ResBlock
scale, shift = self.time_proj(t_emb).chunk(2, dim=1)
h = self.act(self.norm2(h) * (1 + scale[...,None,None]) + shift[...,None,None])
\end{lstlisting}
\highlight{The three pillars --- sinusoidal time embedding, AdaGN conditioning, bottleneck attention ---
reappear in notebook 04 and (with FiLM instead of AdaGN) in notebook 05. Ask: \emph{why does the
network need the timestep at all, if it only ever sees $\mathbf{x}_t$?}}

\subsection{Full U-Net (CODE \#8)}
\whatdoes{Assembles the encoder ($28\to14\to7$ via MaxPool, channels $64\to128\to256$), the
attention bottleneck, and the decoder with ConvTranspose upsampling and concatenated skip connections,
ending in a $1\times1$ conv that outputs the predicted noise. Runs a shape sanity check and prints the
parameter count.}
\begin{saythis}
This is the classic U-shape: contract down to a small spatial bottleneck, then expand back up, with skip
connections carrying fine detail across. The crucial property: input shape equals output shape ---
twenty-eight by twenty-eight in, twenty-eight by twenty-eight out --- because the output \emph{is} the
predicted noise, one value per pixel. The sanity check confirms exactly that.
\end{saythis}
\highlight{Output shape $=$ input shape is non-negotiable for $\varepsilon_\theta$. Skip connections are
what make U-Nets so good at preserving detail.}

\subsection{Training loop (CODE \#9) --- LAUNCH THIS, then go to NB 04}
\whatdoes{The $\mathcal{L}_{\text{simple}}$ training loop. Each step: sample a batch $\mathbf{x}_0$,
a random $t$, fresh noise $\boldsymbol\varepsilon$; form $\mathbf{x}_t$; predict noise with the U-Net;
minimize \texttt{mse\_loss(pred\_noise, noise)}; clip gradients to 1.0; Adam step; cosine LR schedule.
Saves a checkpoint every 5 epochs.}
\begin{saythis}
And here is the entire learning objective, and it is shockingly simple. For each image, pick a random
noise level, actually add that noise, and train the network to predict the noise back out. That is it
--- mean squared error between true and predicted noise. No adversarial game, no likelihood gymnastics,
just regression. I am going to start this cell \emph{now}. It will run for the next fifteen to thirty
minutes. While it trains, we switch to notebook four and look at the different network architectures we
could have plugged in here.
\end{saythis}
\begin{lstlisting}
t     = torch.randint(0, T_STEPS, (x0.size(0),), device=DEVICE)
noise = torch.randn_like(x0)
xt, _ = q_sample(x0, t, noise)
loss  = F.mse_loss(model(xt, t), noise)   # L_simple
\end{lstlisting}
\highlight{This is $\mathcal{L}_{\text{simple}} = \mathbb{E}\big[\|\varepsilon - \varepsilon_\theta(\mathbf{x}_t,t)\|^2\big]$.
Gradient clipping stabilizes early training. \textbf{Start the cell before lecturing.} Ask:
\emph{why is regressing on noise easier to optimize than regressing on the clean image directly?}}

\subsection{Loss curve (CODE \#10)}
\whatdoes{Plots the per-epoch average MSE loss; prints the final loss.}
\begin{saythis}
When we come back, this is the first thing we check: a loss that drops and flattens. The absolute number
is not meaningful --- it is the \emph{shape} that tells us learning happened.
\end{saythis}
\highlight{Noise-prediction MSE plateaus quickly; further gains are mostly in sample quality, not loss.}

\subsection{DDPM sampler (CODE \#11)}
\whatdoes{Implements the ancestral DDPM sampler: start from $\mathbf{x}_T\sim\mathcal{N}(0,I)$, loop
$t=999\ldots0$, predict noise, compute the posterior mean, add posterior-variance noise when $t>0$.
Generates 16 samples. Runs 1000 network evaluations per sample.}
\begin{saythis}
Now the network is the denoiser --- no more oracle. We start from pure noise and run the reverse process
one thousand times. Notice the cost: one thousand forward passes through the U-Net to make a single
batch of digits. That is the DDPM bottleneck, and it is exactly why DDIM and friends exist.
\end{saythis}
\highlight{1000 network calls per sample $=$ slow. Connect forward to the speed problem we solve in
notebook 05 with DDIM. Ask: \emph{how could we cut this to 50 steps without retraining?}}

\subsection{Generated grid, denoising trajectory, real-vs-fake (CODE \#12--14)}
\whatdoes{CODE \#12 shows a $4\times4$ grid of generated digits and saves it. CODE \#13 saves and plots
the denoising trajectory of a single sample from noise to digit. CODE \#14 puts real MNIST (top) next to
generated (bottom).}
\begin{saythis}
Here is the payoff. After only twenty epochs these are clearly digits --- maybe a bit soft, and a few
ambiguous ones, which is honest for twenty epochs. The trajectory plot is my favourite: you literally
watch a digit condense out of static. And side by side with real MNIST, our generations hold up well.
More epochs would sharpen them further.
\end{saythis}
\highlight{Manage expectations: 20 epochs gives good-but-soft digits. The denoising trajectory is the
single best visual for ``what a diffusion model does''. We have now built DDPM end-to-end from scratch.}

% ===========================================================================
\section{Notebook 04 --- Architectures for Diffusion Models}
% ===========================================================================

\textbf{Learning goal.} Drive home that the diffusion process is \emph{architecture-agnostic}: any
network mapping $(\mathbf{x}_t, t)\mapsto$ a same-shape tensor works as $\varepsilon_\theta$. We build
four backbones --- a 2D U-Net, a 2D U-Net with multi-scale attention, a 1D U-Net for sequences, and a
DiT Transformer --- run a forward pass through each, and benchmark parameter counts and speed. The 1D
U-Net is the backbone of Diffusion Policy. \textbf{Cover this entire notebook while NB 03 trains.}

\subsection{Setup, imports, shared primitives (CODE \#1--3)}
\whatdoes{CODE \#1--2 check CUDA / import. CODE \#3 defines the shared building blocks reused by all four
models: \texttt{SinusoidalEmbedding}, \texttt{TimeEmbedMLP}, and an \texttt{AdaGN} module that conditions
activations on the time embedding via scale and shift --- generalized to any number of spatial dimensions.}
\begin{saythis}
Before the four architectures, notice they all share the same time-conditioning machinery from notebook
three: embed the timestep, push it through a small MLP, and inject it as a scale-and-shift. Writing it
once here makes the point that the conditioning mechanism is shared; only the backbone differs.
\end{saythis}
\highlight{The AdaGN here is written dimension-agnostically (works for 1D and 2D) --- foreshadowing the
1D U-Net.}

\subsection{2D U-Net (CODE \#4)}
\whatdoes{Builds \texttt{UNet2D}: residual blocks with AdaGN, self-attention at the $7\times7$ bottleneck
only, MaxPool down / ConvTranspose up with skip connections; runs a forward pass on a
$(2,1,28,28)$ input and prints the parameter count ($\sim$3M).}
\begin{saythis}
This is the DDPM workhorse, essentially the model we just trained in notebook three, re-expressed with
the shared blocks. Attention only at the bottleneck, where the spatial grid is small and global
reasoning is cheap. Around three million parameters.
\end{saythis}
\highlight{Bottleneck-only attention keeps cost down on small images. Ask: \emph{at $7\times7$ the
attention matrix is $49\times49$ --- why would full-resolution attention at $28\times28$ be much more
expensive?}}

\subsection{2D U-Net with multi-scale attention (CODE \#5)}
\whatdoes{Builds \texttt{UNet2DAttentionFull} (Improved DDPM): adds self-attention at \emph{every}
resolution, not just the bottleneck. Forward pass and parameter count ($\sim$5M).}
\begin{saythis}
The improvement from Nichol and Dhariwal: put attention at \emph{every} scale, so the model gets global
context even at high resolution, not only at the bottleneck. It costs more parameters and compute --- the
attention is order $N$-squared in the number of pixels --- but it helps on larger, more detailed images.
\end{saythis}
\highlight{More attention $=$ more global context but $O(N^2)$ cost. This is the central scaling tension
for image diffusion.}

\subsection{1D U-Net for sequences (CODE \#6) --- the Diffusion Policy backbone}
\whatdoes{Builds \texttt{UNet1D}: identical U-Net structure but every \texttt{Conv2d} becomes
\texttt{Conv1d}, operating on \texttt{(B, D, T)} = batch $\times$ feature dim $\times$ sequence length,
with kernel size 5. Forward pass on a $(2,2,16)$ trajectory and parameter count.}
\begin{saythis}
This is the one that matters for the rest of the day. Take the exact same U-Net and swap every
two-dimensional convolution for a one-dimensional one. Now it slides along \emph{time} instead of across
\emph{space}. The input is batch by action-dimension by horizon. This is the backbone of Diffusion
Policy --- the network we will condition on observations and train in notebook five. Same diffusion
maths, same U-Net skeleton, different convolution.
\end{saythis}
\begin{lstlisting}
# Trajectories: (B, D_action, T_horizon); 2D conv -> 1D conv
class UNet1D(nn.Module):
    def __init__(self, in_ch=2, dim=64, t_dim=256):
        self.inc = nn.Conv1d(in_ch, dim, 3, padding=1)  # ...same U-Net, 1D
\end{lstlisting}
\highlight{This is the literal bridge from images to robots: \texttt{Conv2d}$\to$\texttt{Conv1d}. Ask:
\emph{what does a 1D convolution along the time axis capture about a trajectory?} (local temporal
smoothness / short action motifs).}

\subsection{DiT --- Diffusion Transformer (CODE \#7)}
\whatdoes{Builds a simplified DiT (Peebles \& Xie): patchify the image into tokens, linear-project, add a
positional embedding, run $L=6$ Transformer blocks with \emph{AdaLayerNorm} time conditioning, then
unpatchify back to the noise prediction. Forward pass and parameter count ($\sim$15M).}
\begin{saythis}
The fourth option throws away the U-Net entirely and uses a pure Transformer on image patches, with time
injected through adaptive layer norm instead of adaptive group norm. Its advantage is that it has no
built-in locality bias --- it can attend globally from the very first layer --- and it scales beautifully
with data and compute. The cost is that it has many more parameters and is data-hungry.
\end{saythis}
\highlight{DiT scales with compute where U-Nets saturate; but it needs more data. Same diffusion process,
radically different inductive bias.}

\subsection{Benchmark and schematics (CODE \#8--9)}
\whatdoes{CODE \#8 warms up and times 20 forward passes of each model (batch 2) and bar-charts parameter
count and forward-pass time. CODE \#9 draws schematic block diagrams of all four architectures.}
\begin{saythis}
Let's compare them head to head. Read off two things: the DiT has by far the most parameters, and the
1D U-Net is the leanest because sequences are low-dimensional. That last point is exactly why Diffusion
Policy can run fast enough to control a robot. All four, remember, are interchangeable as the noise
predictor --- the choice is about your data, not about the diffusion maths.
\end{saythis}
\highlight{The 1D U-Net's small footprint is what makes real-time trajectory diffusion practical ---
the perfect segue to notebook 05. By now, notebook 03 should be finished: go back and show its samples.}

% ===========================================================================
\section{Notebook 05 --- Diffusion Policy on the Real Push-T Dataset}
% ===========================================================================

\textbf{Learning goal.} Train a real \emph{Diffusion Policy} (Chi et al., 2023) on the Push-T
benchmark, where an end-effector pushes a T-shaped block to a target. We engineer observations and
delta-actions, normalize them, build a conditional 1D U-Net that denoises action chunks given the past
observations, train it, and compare DDPM versus DDIM sampling --- showing why DDIM's $\sim$10$\times$
speed-up is what makes 10~Hz robot control feasible. \textbf{Strategy: run setup + download + model,
launch training (CODE \#12) early, lecture, then show results.}

\subsection{Setup with extra deps (CODE \#1)}
\whatdoes{Beyond the usual CUDA check, this Setup cell \texttt{pip install}s \texttt{zarr} and
\texttt{gdown} if \texttt{zarr} is not importable --- needed to read and download the dataset.}
\begin{saythis}
This notebook needs two extra libraries that Colab does not ship: \texttt{zarr}, to read the dataset's
chunked array format, and \texttt{gdown}, to fetch it from Google Drive. The setup cell installs them
if they are missing. Run it first.
\end{saythis}
\highlight{Note: the notebook reads the dataset with zarr v3 using \texttt{zarr\_format=2}, which
requires Python $\ge 3.11$ --- fine on current Colab. If you ever run this on an older local Python, this
is the first thing that breaks.}

\subsection{Imports and configuration (CODE \#2--3)}
\whatdoes{CODE \#2 imports and seeds. CODE \#3 is the config: \texttt{D\_OBS = 4}, \texttt{D\_ACTION = 2},
\textbf{\texttt{T\_OBS = 2}} (observation horizon), \textbf{\texttt{T\_PRED = 16}} (prediction horizon),
\texttt{T\_STEPS = 100} diffusion steps, \texttt{BATCH\_SIZE = 64}, \textbf{\texttt{NUM\_EPOCHS = 100}},
\texttt{LR = 1e-3}.}
\begin{saythis}
These five numbers define the policy. We look at the last two observations --- that is the observation
horizon, $T_o$ equals two --- and from them we predict the next sixteen actions in one shot, the
prediction horizon $T_p$ equals sixteen. Predicting a whole \emph{chunk} of actions, not one step, is a
core Diffusion Policy idea: it gives temporally consistent, non-jittery motion. Also notice we use only
one hundred diffusion steps here, not a thousand --- action chunks are small and easy, so far fewer
steps suffice.
\end{saythis}
\highlight{$T_o=2$, $T_p=16$, $D_a=2$, $T=100$. Action chunking $=$ temporal consistency. One hundred
epochs is $\approx$30--40~min on a T4 --- this is the cell we will background. Ask: \emph{why predict 16
actions at once instead of replanning every single step?}}

\subsection{Dataset download (CODE \#4) --- the added gdown cell}
\whatdoes{If the dataset directory is absent, this cell downloads \texttt{pusht\_cchi\_v7\_replay.zarr.zip}
($\approx$31~MB) from Google Drive via \texttt{gdown} and unzips it into \texttt{data/pusht/}; otherwise it
reports the dataset is already present.}
\begin{saythis}
This cell fetches the real dataset --- about thirty-one megabytes --- straight from Google Drive and
unzips it. It only does this once; on a re-run it sees the folder and skips. Give it a few seconds. This
is real robot demonstration data, not the synthetic toys from notebook one: two hundred and six expert
episodes of pushing a T-block.
\end{saythis}
\begin{lstlisting}
if not os.path.isdir(os.path.join(DATA_DIR, "data")):
    gdown.download(id="1KY1InLurp...", output=_zip, quiet=False)
    with zipfile.ZipFile(_zip) as zf:
        zf.extractall(DATA_DIR)
\end{lstlisting}
\highlight{Common failure modes: no internet, a Google Drive quota error, or a stale partial zip. If the
download fails, the fix is to delete \texttt{data/pusht\_cchi\_v7\_replay.zarr.zip} and re-run. Watch the
\texttt{gdown} progress bar before moving on.}

\subsection{Load the dataset (CODE \#5)}
\whatdoes{Opens the zarr groups with \texttt{zarr\_format=2}, reads \texttt{data/state} $(N,5)$,
\texttt{data/action} $(N,2)$ and \texttt{meta/episode\_ends} $(206,)$; derives per-episode starts and
lengths and prints shapes, ranges, and episode statistics ($N\approx25{,}650$ timesteps, 206 episodes).}
\begin{saythis}
Here is the structure of the data. The state is five numbers per timestep: the agent's x and y, the
block's x and y, and the block's angle. The action is the absolute target position the controller was
commanded. And \texttt{episode\_ends} tells us where one demonstration stops and the next begins ---
essential, because we must never let a training window straddle two episodes.
\end{saythis}
\highlight{Positions are in pixel space (0--512). Episode boundaries matter for windowing. Ask:
\emph{why would mixing two episodes in one observation-action window corrupt training?}}

\subsection{Feature engineering (CODE \#6)}
\whatdoes{Computes per-episode velocity as finite differences (zero at each episode start), builds the
observation \texttt{[agent\_x, agent\_y, vx, vy]} $(N,4)$, and the \emph{delta} action
\texttt{action - agent\_pos} $=$ \texttt{[dx, dy]} $(N,2)$. Prints statistics.}
\begin{saythis}
Two design decisions here. First, we enrich the observation with velocity --- position alone does not
tell you which way the agent is moving, and that matters for prediction. We compute velocity as the
difference between consecutive positions, reset to zero at each episode start. Second, and importantly,
we predict a \emph{delta} action: not ``go to pixel 300,'' but ``move 40 pixels right from where you
are.'' Relative targets are far easier to learn and generalize.
\end{saythis}
\highlight{Delta (relative) actions $=$ better generalization than absolute targets. Velocity makes the
observation Markov-ish. Ask: \emph{why reset velocity to zero at episode boundaries?}}

\subsection{Normalization statistics (CODE \#7)}
\whatdoes{Computes per-dimension mean and std of observations and actions over the whole dataset (std
clipped at $10^{-3}$); prints them. These are stored and later used to normalize inputs and
\emph{de}-normalize generated actions.}
\begin{saythis}
Same lesson as notebook one, now for real: standardize every channel to zero mean, unit variance. We
compute these statistics once over the whole dataset and keep them. They get saved into every checkpoint
too, because at deployment we must apply the \emph{exact same} normalization and then invert it on the
output --- otherwise the predicted motions are garbage.
\end{saythis}
\highlight{The stats are saved \emph{in the checkpoint} --- normalization is part of the policy, not an
afterthought. Forgetting to de-normalize at inference is the classic bug.}

\subsection{Dataset visualization (CODE \#8)}
\whatdoes{Four-panel figure: agent trajectories colored by episode, agent-vs-block position density, the
delta-action $[dx,dy]$ scatter with 1$\sigma$/2$\sigma$ ellipses, and the episode-length histogram with
CDF; plus a printed summary table.}
\begin{saythis}
Spend a moment on the bottom-left panel: the distribution of delta-actions. It is roughly centred on
zero, slightly spread --- and crucially, it is \emph{multi-modal} in places: from a given state there are
often several reasonable pushes. That multi-modality is precisely why a diffusion policy beats a simple
regression-to-the-mean controller, which would average two good options into one bad one.
\end{saythis}
\highlight{Multi-modality of demonstrations is the core motivation for generative (vs regression)
policies. The action ellipses show the scale the network must learn to reproduce.}

\subsection{Dataset class --- observation/action windows (CODE \#9)}
\whatdoes{Defines \texttt{PushTDataset}: normalizes obs and actions, then enumerates valid index pairs so
that each sample is the last $T_o$ observations \texttt{(T\_obs, D\_obs)} and the next $T_p$ actions
\texttt{(T\_pred, D\_action)}, excluding windows that cross an episode boundary. Builds the
\texttt{DataLoader} and prints shapes.}
\begin{saythis}
This turns the long demonstration streams into training examples. Each example is a little window: here
are the last two observations, and here are the sixteen actions the expert took next. We slide this
window through every episode but stop it from ever spanning two episodes. The result is tens of
thousands of observation-to-action-chunk pairs --- our supervised dataset.
\end{saythis}
\highlight{The model learns $p(\text{action chunk} \mid \text{recent observations})$. The boundary
exclusion in the index loop is the subtle correctness detail.}

\subsection{Noise schedule (CODE \#10)}
\whatdoes{The same cosine schedule, now with \texttt{T\_STEPS = 100}; precomputes all schedule tensors and
defines \texttt{q\_sample} for action chunks of shape \texttt{(B, D\_action, T\_pred)}.}
\begin{saythis}
The diffusion schedule is identical to the image case --- only the tensor shape changed, from images to
action chunks, and we use a hundred steps instead of a thousand. I want you to feel how little changes:
the entire forward process is the same one line we have seen all day.
\end{saythis}
\highlight{Reinforce the unifying theme: same schedule, same \texttt{q\_sample}, different data shape.}

\subsection{Conditional 1D U-Net (CODE \#11)}
\whatdoes{Defines \texttt{DiffusionPolicy1D}: a 1D U-Net (the notebook-04 backbone) whose residual blocks
are \emph{FiLM}-conditioned on a context vector built by concatenating the sinusoidal time embedding with
an MLP-embedding of the flattened observation window. \texttt{forward(noisy\_action, obs, t)} returns the
predicted noise on the action chunk. Channels $128\to256\to512$. Sanity check and parameter count.}
\begin{saythis}
Here it all comes together. The backbone is the 1D U-Net from notebook four. The new ingredient is
\emph{conditioning}: the network must denoise the action chunk \emph{given} the observation. We do this
with FiLM --- we embed the timestep and the observations, concatenate them into one context vector, and
that context produces a scale-and-shift inside every block, just like AdaGN but driven by the
observation too. So the network's question is: ``given what I just saw, and given this noise level, what
noise did I add to this action chunk?''
\end{saythis}
\begin{lstlisting}
def forward(self, noisy_action, obs, t):
    cond = torch.cat([self.time_emb(t),
                      self.obs_emb(obs.flatten(1))], dim=1)  # FiLM context
    # ... 1D U-Net blocks, each FiLM-conditioned on `cond`
\end{lstlisting}
\highlight{FiLM conditioning $=$ time embedding $\oplus$ observation embedding $\to$ per-block
scale/shift. This is the unconditional DDPM of NB03 plus an observation conditioning signal. Ask:
\emph{what would the policy degenerate into if we removed the observation from the conditioning?}}

\subsection{Training (CODE \#12) --- LAUNCH THIS, then lecture}
\whatdoes{AdamW + cosine LR. Each step: permute the action batch to \texttt{(B, D\_action, T\_pred)}, draw
a random $t$ and noise, form the noisy action chunk, predict noise conditioned on the observation, and
minimize MSE. Gradient clip at 1.0. Saves a checkpoint (with model + normalization stats) every 10
epochs over 100 epochs.}
\begin{saythis}
The training objective is identical to MNIST: predict the noise, minimize mean squared error. The only
difference is that the prediction is now conditioned on the observation. I will start this cell now ---
it runs for about thirty to forty minutes, and saves a checkpoint every ten epochs, so even if we stop
early we have something usable. While it trains, let's talk through the Diffusion Policy idea and how it
bridges to the lecture.
\end{saythis}
\highlight{Same $\mathcal{L}_{\text{simple}}$, now conditional. Checkpoint-every-10 is your safety net.
\textbf{Start the cell, then lecture.} If completely out of time, use CODE \#18 to load a checkpoint.}

\subsection{Loss curve (CODE \#13)}
\whatdoes{Plots the per-epoch MSE on a log scale and prints the final loss.}
\begin{saythis}
On a log scale you can see the loss drop over more than an order of magnitude. As before, the shape ---
steady decrease then plateau --- is what tells us the policy has learned the demonstration distribution.
\end{saythis}

\subsection{DDPM and DDIM samplers (CODE \#14)}
\whatdoes{Defines two inference functions, both taking a raw observation window, normalizing it,
generating $N$ action chunks, and de-normalizing back to pixel-delta space. \texttt{sample\_actions\_ddpm}
runs all 100 stochastic reverse steps; \texttt{sample\_actions\_ddim} runs only \texttt{ddim\_steps=10}
deterministic steps via the predict-$\hat{\mathbf{x}}_0$-then-jump rule.}
\begin{saythis}
Two ways to generate actions from the same trained network. DDPM runs the full hundred-step stochastic
reverse process. DDIM, the fast one, skips most steps: it estimates the clean action chunk and jumps
ahead, in just ten steps. Both end with the de-normalization step --- mapping the network's output back
into real pixel-delta units the robot can execute. Notice the network weights are identical; only the
sampler changes.
\end{saythis}
\begin{lstlisting}
# DDIM step on action chunks: predict x0, jump to previous level
x0_hat = ((x - torch.sqrt(1-ab_i)*eps) / torch.sqrt(ab_i)).clamp(-5,5)
x      = torch.sqrt(ab_p) * x0_hat + torch.sqrt(1-ab_p) * eps
\end{lstlisting}
\highlight{Same weights, two samplers --- the lesson from NB02 made practical. The \texttt{clamp} on
$\hat{\mathbf{x}}_0$ stabilizes the few-step DDIM. De-normalization is mandatory.}

\subsection{Visualize generated action chunks (CODE \#15)}
\whatdoes{Picks three observation windows from episode 0, samples 15 action chunks with each sampler,
integrates each delta-action chunk into an absolute path from the current position, and overlays the
samples (colored) against the ground-truth demonstration (black) and the current position.}
\begin{saythis}
This is the moment of truth. For three different situations, we ask the policy for fifteen candidate
action chunks and draw them. Two things to look for. First, the coloured samples cluster around the
black ground-truth path --- the policy has learned sensible pushes. Second, the samples are \emph{varied},
not identical: that spread \emph{is} the multi-modality we wanted, the policy proposing several
reasonable plans. Compare the DDPM row with the DDIM row --- ten steps gives essentially the same quality
as a hundred.
\end{saythis}
\highlight{Samples should bracket the ground truth and show diversity. DDIM(10) $\approx$ DDPM(100) in
quality. Ask: \emph{if all 15 samples were identical, what would that tell us about the policy?}}

\subsection{DDPM vs DDIM speed (CODE \#16) --- the real-time argument}
\whatdoes{Times one query (10 samples) for DDPM (100 steps) versus DDIM (10 steps), prints both times and
the speed-up, and checks feasibility against a 10~Hz control budget (100~ms per query).}
\begin{saythis}
And here is why we care. DDPM takes the full hundred steps; DDIM takes ten and is roughly an order of
magnitude faster. Now apply the robot's constraint: a ten-hertz controller has just one hundred
milliseconds to decide its next move. The full DDPM sampler is usually too slow to hit that budget; the
ten-step DDIM sampler fits comfortably. \emph{That} is the practical reason DDIM-style fast sampling is
not a luxury but a requirement for closing the loop on a real robot.
\end{saythis}
\highlight{This is the headline result of the session: $\sim$10$\times$ speed-up makes 10~Hz control
feasible without retraining. Train once, sample fast. Ask: \emph{what else could we trade --- horizon,
model size --- to fit the 100~ms budget?}}

\subsection{Action distribution analysis and checkpoint loading (CODE \#17--18)}
\whatdoes{CODE \#17 samples many action chunks across 50 observations (DDIM) and compares the generated
$[dx,dy]$ distribution against the ground-truth distribution side by side. CODE \#18 defines
\texttt{load\_checkpoint} (restoring weights \emph{and} normalization stats) and lists available
checkpoints; the load line is commented out.}
\begin{saythis}
Finally, a population-level check: do the actions the policy generates have the same overall distribution
as the real data? The two scatter clouds should look alike in shape and spread --- if they do, the policy
has matched the demonstration distribution, not just a few examples. And the very last cell shows how to
load any saved checkpoint, including its normalization statistics, so you can skip training next time and
go straight to sampling.
\end{saythis}
\highlight{Distribution match (not just per-query plausibility) is the real success criterion.
\texttt{load\_checkpoint} restores normalization stats too --- reuse it to demo without retraining.}

% ===========================================================================
\section{Closing --- Key Takeaways and the Bridge to the Lecture}
% ===========================================================================

\paragraph{What we proved with our hands today.}
\begin{itemize}
  \item \textbf{One process, many modalities.} The same forward/reverse diffusion maths drove images
  (NB01--03), and trajectories (NB05). Only the data \emph{shape} and the network \emph{backbone}
  changed. We literally turned an image U-Net into a trajectory U-Net by swapping \texttt{Conv2d} for
  \texttt{Conv1d} (NB04).
  \item \textbf{Forward is a one-liner; the network learns to reverse it.} The closed-form
  $\mathbf{x}_t = \sqrt{\bar\alpha_t}\mathbf{x}_0 + \sqrt{1-\bar\alpha_t}\boldsymbol\varepsilon$ makes
  training cheap, and the objective is just MSE noise prediction,
  $\mathcal{L}_{\text{simple}} = \mathbb{E}\|\varepsilon - \varepsilon_\theta(\mathbf{x}_t,t)\|^2$
  (NB02--03, NB05).
  \item \textbf{Normalization is part of the model.} Every modality is scaled to roughly
  $\mathcal{N}(\mathbf{0},\mathbf{I})$; in NB05 the normalization statistics live inside the checkpoint
  and must be inverted at inference.
  \item \textbf{Architecture is a choice, not a constraint.} 2D U-Net, multi-scale-attention U-Net,
  1D U-Net, and DiT are all interchangeable noise predictors (NB04); pick by data and compute.
  \item \textbf{Sampling speed decides real-time feasibility.} DDPM's 100--1000 steps are slow; DDIM
  reaches comparable quality in $\sim$10 steps, giving the $\sim$10$\times$ speed-up that lets a
  Diffusion Policy run inside a 10~Hz control loop (NB05).
\end{itemize}

\paragraph{The bridge to the lecture.} The practical path we walked ---
\textbf{DDPM} (the from-scratch denoiser, NB03) $\rightarrow$ \textbf{DDIM} (the same model, fast
deterministic sampling, NB02 \& NB05) $\rightarrow$ \textbf{architectures} (choosing the backbone for
your modality, NB04) $\rightarrow$ \textbf{Diffusion Policy} (conditional 1D U-Net on action chunks,
NB05) --- mirrors the lecture's arc. The crucial robotics consequence we demonstrated empirically is the
last one: \emph{a Diffusion Policy is trained once but must be sampled in real time}, and that is exactly
what DDIM (and, beyond this session, DPM-Solver and consistency models) buys us. Notebook 05 made the
$\sim$10$\times$ DDIM speed-up concrete and tied it to the 100~ms-per-decision budget of a 10~Hz
controller --- the practical link between the generative-model theory and deploying a learned policy on a
compliant robot.

\paragraph{Where to go next.} Natural extensions raised by the notebooks: richer observations (include
the block state, longer observation windows, or visual inputs via a CNN encoder), a Transformer backbone
for longer horizons, classifier-free guidance and latent diffusion for images, and faster solvers for
even tighter control loops.

\end{document}
